\chapter{The Relativistic String: Nambu--Goto and Polyakov}\label{ch:classical}

What is the simplest relativistic dynamics that a one-dimensional object can have, and how
does one solve it? This chapter answers that question for the classical closed string. The
answer has three layers. First, an action: the area of the surface that the string sweeps
out in spacetime, in the same way that the action of a free particle is the length of its
worldline. Second, a reformulation, the Polyakov action, which removes the square root at the
price of a new field and a new gauge symmetry. Third, a solution: in a suitable gauge the
string obeys the free wave equation, its general motion is a superposition of left-moving and
right-moving waves, and what remains of the gauge symmetry appears as an infinite set of
quadratic constraints on the Fourier modes. A student should care because every later
chapter starts from the end point of this one. The momentum and winding of
\cref{ch:circle}, the duality $R\mapsto\ap/R$ of \cref{ch:tduality}, the Narain lattice of
\cref{ch:narain} are all statements about the zero modes and the constraints $L_0=\tilde
L_0=0$ derived here. The chapter also marks a boundary of this book's formal content: almost
nothing in it is machine-checked, and \cref{sec:classical:lean} explains why that is the
expected state of affairs and not an omission.

Conventions are those of the whole book: $\hbar=c=1$, mostly-plus signature
$\eta_{\mu\nu}=\diag(-1,+1,\dots,+1)$ in $D$ spacetime dimensions
\source{Tong §1.1, ll.~688--693}, and $A\cdot B=A^\mu B^\nu\eta_{\mu\nu}$, $A^2=A\cdot A$. A
timelike vector has $A^2<0$.

%=====================================================================================
\section{Warm-up: the relativistic point particle}\label{sec:classical:particle}

Every feature of the string action has a simpler counterpart for the particle. It pays to
see each one there first.

\subsection{From a frame-dependent to a covariant action}

Fix an inertial frame with coordinates $X^\mu=(t,\vec x)$. A free particle of mass $m$ has the
action
\begin{equation}\label{eq:classical:particle-frame}
  S=-m\int \dd t\,\sqrt{1-\dot{\vec x}\cdot\dot{\vec x}}\,,\qquad \dot{\vec x}=\dd\vec x/\dd t
\end{equation}
\source{Tong §1.1, ll.~694--712}. One checks it by computing the conjugate momentum and the
energy. With $L=-m\sqrt{1-\dot{\vec x}^{\,2}}$,
\begin{equation}\label{eq:classical:particle-p}
  \vec p=\frac{\partial L}{\partial\dot{\vec x}}=\frac{m\dot{\vec x}}{\sqrt{1-\dot{\vec x}^{\,2}}},
  \qquad
  E=\vec p\cdot\dot{\vec x}-L=\frac{m\dot{\vec x}^{\,2}}{\sqrt{1-\dot{\vec x}^{\,2}}}+m\sqrt{1-\dot{\vec x}^{\,2}}
   =\frac{m}{\sqrt{1-\dot{\vec x}^{\,2}}}\,,
\end{equation}
so that $E^2-\vec p^{\,2}=m^2(1-\dot{\vec x}^{\,2})/(1-\dot{\vec x}^{\,2})=m^2$, the relativistic
dispersion relation $E=\sqrt{m^2+\vec p^{\,2}}$.

The action \eqref{eq:classical:particle-frame} is correct but treats $t$ and $\vec x$
differently: $\vec x$ is a dynamical variable, $t$ a label. Lorentz transformations mix them, so
Lorentz invariance is hidden. The remedy is to let $t$ become a dynamical variable too, and
to pay for it with a gauge symmetry \source{Tong §1.1, ll.~713--738}. Introduce an
arbitrary parameter $\tau$ along the worldline\index{worldline} and write
\begin{equation}\label{eq:classical:particle-cov}
  S=-m\int\dd\tau\,\sqrt{-\dot X^\mu\dot X^\nu\eta_{\mu\nu}}\,,\qquad \dot X^\mu=\dd X^\mu/\dd\tau,
  \qquad \mu=0,\dots,D-1
\end{equation}
\source{Tong §1.1, ll.~739--760}. Since $\dd s=\sqrt{-\dd X\cdot\dd X}$ is the proper time
elapsed along an element of a timelike curve, the action is $-m$ times the total proper
time, that is, $-m$ times the Lorentzian length of the worldline. It is manifestly invariant
under the Poincar\'e transformations\index{Poincar\'e invariance}
$X^\mu\to\Lambda^\mu{}_\nu X^\nu+c^\mu$, because only the invariant $\dot X^2$ enters
\source{Tong §1.1, ll.~815--822}.

\subsection{Reparametrization invariance is a gauge symmetry}

The action \eqref{eq:classical:particle-cov} now has $D$ functions $X^\mu(\tau)$, one more than
the $D-1$ functions $\vec x(t)$ we started with. The extra one is not physical, because the
action does not depend on the choice of parameter. Let $\tilde\tau(\tau)$ be any monotonic
function. Then $\dd\tau=\dd\tilde\tau\,|\dd\tau/\dd\tilde\tau|$ and
$\dd X^\mu/\dd\tau=(\dd X^\mu/\dd\tilde\tau)(\dd\tilde\tau/\dd\tau)$. The square root is
homogeneous of degree one in the velocities, so the two Jacobian factors cancel:
\begin{equation}\label{eq:classical:reparam}
  S=-m\int\dd\tilde\tau\,\sqrt{-\frac{\dd X^\mu}{\dd\tilde\tau}\frac{\dd X^\nu}{\dd\tilde\tau}\eta_{\mu\nu}}
\end{equation}
\source{Tong §1.1, ll.~765--782}. This \emph{reparametrization invariance}%
\index{reparametrization invariance} is a redundancy of the description, not a symmetry
relating different physical states. We can use it to choose $\tau=X^0(\tau)=t$; then
$-\dot X^2=1-\dot{\vec x}^{\,2}$ and \eqref{eq:classical:particle-cov} reduces to
\eqref{eq:classical:particle-frame} \source{Tong §1.1, ll.~783--797}.

The redundancy is also visible in the momenta. From \eqref{eq:classical:particle-cov},
\begin{equation}\label{eq:classical:particle-momenta}
  p_\mu=\frac{\partial L}{\partial\dot X^\mu}=\frac{m\,\dot X^\nu\eta_{\mu\nu}}{\sqrt{-\dot X^2}}
  \qquad\Longrightarrow\qquad
  p_\mu p^\mu=\frac{m^2\dot X^2}{-\dot X^2}=-m^2 .
\end{equation}
The $D$ momenta are not independent; they satisfy the \emph{mass-shell constraint}%
\index{mass-shell constraint} $p^2+m^2=0$ identically, for every trajectory, before any
equation of motion is used \source{Tong §1.1, ll.~798--814}. A second consequence of the same
homogeneity: the canonical Hamiltonian vanishes,
\begin{equation}\label{eq:classical:Hzero}
  H=\dot X^\mu p_\mu-L=\frac{m\dot X^2}{\sqrt{-\dot X^2}}+m\sqrt{-\dot X^2}=-m\sqrt{-\dot X^2}+m\sqrt{-\dot X^2}=0
\end{equation}
\source{Tong §1.1.1, ll.~826--840}. There is no evolution in $\tau$ because $\tau$ means
nothing; all the dynamics sits in the constraint. In the quantum theory the constraint is
imposed on wavefunctions, $(p^2+m^2)\Psi=0$ with $p_\mu=-\ii\partial/\partial X^\mu$, which is
the Klein--Gordon equation \source{Tong §1.1.1, ll.~848--857}.

\begin{physicsbox}[Physical picture: a constraint for every gauge symmetry]
A gauge symmetry with one arbitrary function of $\tau$ removes one function from the
physical data and produces one constraint on the momenta. Remember this bookkeeping. The
string has worldsheet reparametrizations with \emph{two} arbitrary functions of \emph{two}
variables; it will have two constraints at every point of the string, and their Fourier
modes will be the $L_n$ and $\tilde L_n$ of \cref{sec:classical:modes}.
\end{physicsbox}

\subsection{The einbein: removing the square root}\label{sec:classical:einbein}

There is a second action for the particle. Introduce an additional worldline field $e(\tau)>0$,
the \emph{einbein}\index{einbein}, and set
\begin{equation}\label{eq:classical:einbein}
  S=\frac12\int\dd\tau\,\Big(e^{-1}\dot X^2-e\,m^2\Big)
\end{equation}
\source{Tong §1.1.2, ll.~885--897}. No derivative of $e$ appears, so its equation of motion is
algebraic: $\partial L/\partial e=-\tfrac12 e^{-2}\dot X^2-\tfrac12 m^2=0$, that is
\begin{equation}\label{eq:classical:einbein-eom}
  \dot X^2+e^2m^2=0,\qquad e=\frac{\sqrt{-\dot X^2}}{m}.
\end{equation}
Substituting back,
$\tfrac12\big(e^{-1}\dot X^2-e\,m^2\big)=\tfrac12\big(-m\sqrt{-\dot X^2}-m\sqrt{-\dot X^2}\big)
=-m\sqrt{-\dot X^2}$, which is \eqref{eq:classical:particle-cov}
\source{Tong §1.1.2, ll.~913--918}. The two actions are classically equivalent. The new one
has two advantages: it makes sense for $m=0$, and it is quadratic in $X$, which is what a path
integral needs \source{Tong §1.1.2, ll.~919--922}. Writing $e=\sqrt{-g_{\tau\tau}}$ shows
that $e$ is a metric on the one-dimensional worldline, so \eqref{eq:classical:einbein} is
``$D$ scalar fields coupled to one-dimensional gravity'' \source{Tong §1.1.2, ll.~898--912}.
Under $\tau\to\tau-\eta(\tau)$ the fields transform as $\delta X^\mu=\eta\,\dot X^\mu$ (a
scalar) and $\delta e=\frac{\dd}{\dd\tau}(\eta e)$ (a density)
\source{Tong §1.1.2, ll.~923--938}.

The step from \eqref{eq:classical:particle-cov} to \eqref{eq:classical:einbein} is the model
for the step from the Nambu--Goto action to the Polyakov action below.

%=====================================================================================
\section{The Nambu--Goto action}\label{sec:classical:ng}

\subsection{Worldsheet, induced metric, area}

A particle sweeps out a worldline; a string sweeps out a two-dimensional surface, the
\emph{worldsheet}\index{worldsheet}. Parametrize it by a timelike coordinate $\tau$ and a
spacelike coordinate $\sigma$, collected as $\sigma^\alpha=(\tau,\sigma)$, $\alpha=0,1$. The
string is described by $D$ functions $X^\mu(\tau,\sigma)$, a map from the worldsheet into
spacetime, which in this context is called the \emph{target space}\index{target space}. This
chapter treats closed strings: $\sigma$ is periodic,
\begin{equation}\label{eq:classical:periodicity}
  \sigma\in[0,2\pi),\qquad X^\mu(\tau,\sigma+2\pi)=X^\mu(\tau,\sigma)
\end{equation}
\source{Tong §1.2, ll.~945--968}; see \cref{fig:classical:worldsheet}. Open strings appear in
\cref{ch:dbranes}.

\begin{figure}[t]
\centering
\begin{tikzpicture}[>=Stealth,scale=0.95]
  % axes
  \draw[->] (-0.3,0)--(2.3,0) node[below]{$\vec x$};
  \draw[->] (0,-0.3)--(0,4.3) node[left]{$X^0$};
  % worldline
  \draw[thick] (0.9,0.2) .. controls (1.4,1.4) and (0.6,2.6) .. (1.3,4.0);
  \fill (1.04,1.2) circle (1.3pt) node[right]{\small $\tau_1$};
  \fill (1.0,2.6) circle (1.3pt) node[right]{\small $\tau_2$};
  \node at (1.0,-0.7){\small (a) worldline $X^\mu(\tau)$};
  % worldsheet
  \begin{scope}[xshift=5.5cm]
    \draw[->] (-1.9,0)--(2.4,0) node[below]{$\vec x$};
    \draw[->] (-1.6,-0.3)--(-1.6,4.3) node[left]{$X^0$};
    \draw[thick] (-0.9,0.5) .. controls (-0.5,1.6) and (-1.1,2.8) .. (-0.6,3.9);
    \draw[thick] (1.1,0.5) .. controls (1.6,1.6) and (0.9,2.8) .. (1.5,3.9);
    \draw[thick] (0.1,0.5) ellipse (1.0 and 0.22);
    \draw[thick] (0.45,3.9) ellipse (1.05 and 0.22);
    \draw[densely dashed] (-0.76,1.6) arc (180:360:0.98 and 0.2);
    \draw[densely dashed] (-0.88,2.7) arc (180:360:1.0 and 0.2);
    \draw[->] (1.9,1.4)--(1.9,2.6) node[right]{$\tau$};
    \draw[->] (0.0,4.3) .. controls (0.3,4.42) and (0.7,4.42) .. (1.0,4.3) node[right]{$\sigma$};
    \node at (0.3,-0.7){\small (b) worldsheet $X^\mu(\tau,\sigma)$};
  \end{scope}
\end{tikzpicture}
\caption{(a) A particle traces a worldline; the parameter $\tau$ labels its points and has no
physical meaning. (b) A closed string traces a tube, the worldsheet. Dashed curves are lines
of constant $\tau$, one snapshot of the string each; $\sigma\sim\sigma+2\pi$ runs around the
tube.}
\label{fig:classical:worldsheet}
\end{figure}

We want an action that depends on the surface and not on the coordinates $\sigma^\alpha$ used
to describe it. For the particle the answer was the length of the worldline. For the string it
is the area of the worldsheet. To compute it, note that the flat metric of spacetime induces a
metric on the surface: two nearby worldsheet points separated by $\dd\sigma^\alpha$ are
separated in spacetime by $\dd X^\mu=\partial_\alpha X^\mu\dd\sigma^\alpha$, hence by the
invariant interval $\dd s^2=\eta_{\mu\nu}\partial_\alpha X^\mu\partial_\beta X^\nu\,
\dd\sigma^\alpha\dd\sigma^\beta$. The \emph{induced metric}\index{induced metric} (or pull-back
of $\eta$) is therefore
\begin{equation}\label{eq:classical:induced}
  \gamma_{\alpha\beta}=\frac{\partial X^\mu}{\partial\sigma^\alpha}\frac{\partial X^\nu}{\partial\sigma^\beta}\,\eta_{\mu\nu}
  =\begin{pmatrix}\dot X^2 & \dot X\cdot X'\\ \dot X\cdot X' & X'^2\end{pmatrix},
  \qquad \dot X^\mu=\partial_\tau X^\mu,\quad X'^\mu=\partial_\sigma X^\mu,
\end{equation}
and the \emph{Nambu--Goto action}\index{Nambu--Goto action} is
\begin{equation}\label{eq:classical:ng}
  S_{\rm NG}=-T\int\dd^2\sigma\,\sqrt{-\det\gamma}
   =-T\int\dd^2\sigma\,\sqrt{(\dot X\cdot X')^2-\dot X^2\,X'^2}
\end{equation}
\source{Tong §1.2, ll.~969--1016}. The second form follows from $\det\gamma=\dot X^2X'^2-(\dot
X\cdot X')^2$. For a physical worldsheet $\dot X$ is timelike and $X'$ spacelike, so
$\det\gamma<0$ and the square root is real. The constant $T$ is fixed below.

Why is $\sqrt{|\det\gamma|}\,\dd^2\sigma$ an area? Consider first a surface in Euclidean space,
$\vec X(\tau,\sigma)$. The coordinate cell $\dd\tau\,\dd\sigma$ is mapped to the parallelogram
spanned by $\vec u=\partial_\tau\vec X\,\dd\tau$ and $\vec v=\partial_\sigma\vec X\,\dd\sigma$, with
angle $\theta$ between them. Its area is
\begin{equation}\label{eq:classical:area}
  |\vec u||\vec v|\sin\theta=\sqrt{\vec u^{\,2}\vec v^{\,2}(1-\cos^2\theta)}
   =\sqrt{\vec u^{\,2}\,\vec v^{\,2}-(\vec u\cdot\vec v)^2}
   =\sqrt{\det\gamma}\;\dd\tau\,\dd\sigma
\end{equation}
\source{Tong §1.2, ll.~1018--1060}. The quantity under the root is non-negative by the
Cauchy--Schwarz inequality. In Minkowski signature the sign under the root flips, which is the
origin of the minus sign in \eqref{eq:classical:ng}. The area is also manifestly independent of
coordinates: under $\sigma^\alpha\to\tilde\sigma^\alpha(\sigma)$ with Jacobian matrix
$J^\gamma{}_\alpha=\partial\sigma^\gamma/\partial\tilde\sigma^\alpha$, the induced metric
becomes $\tilde\gamma=J^{\mathsf T}\gamma J$, so $\sqrt{-\det\tilde\gamma}=|\det
J|\sqrt{-\det\gamma}$, while $\dd^2\tilde\sigma=|\det J|^{-1}\dd^2\sigma$. This is the
two-dimensional version of the cancellation that led to \eqref{eq:classical:reparam}.

\subsection{Tension, \texorpdfstring{$\ap$}{alpha'} and the string length}

To interpret $T$, take a string momentarily at rest. Choose the parametrization so that
$X^0=t=R\tau$ with $R$ a constant of dimension length, and suppose $\dot{\vec x}=0$ at the
instant considered. Then $\dot X^2=-R^2$, $\dot X\cdot X'=0$, $X'^2=(\dd\vec x/\dd\sigma)^2$, and
\begin{equation}\label{eq:classical:tension}
  S_{\rm NG}=-T\int\dd\tau\,\dd\sigma\,R\sqrt{(\dd\vec x/\dd\sigma)^2}
   =-T\int\dd t\,\times(\text{length of the string})
\end{equation}
\source{Tong §1.2, ll.~1061--1086}. For a system at rest the Lagrangian is minus the potential
energy, hence
\[
  \text{potential energy}=T\times\text{length}.
\]
$T$ is the energy per unit length, the \emph{tension}\index{string tension}. The energy
grows linearly with the length, unlike a Hooke spring, whose energy grows quadratically; a
classical string lowers its energy by shrinking \source{Tong §1.2, ll.~1087--1096}.

The book-wide parametrization of the tension is
\begin{equation}\label{eq:classical:alphaprime}
  T=\frac{1}{2\pi\ap},\qquad \ap=\ell_s^2
\end{equation}
\source{Tong §1.2, ll.~1097--1123}, where $\ap$\index{alpha-prime@$\ap$} is called the
universal Regge slope for historical reasons and $\ell_s$\index{string length} is the string
length. Dimensional analysis confirms the powers. With $\hbar=1$ the action is dimensionless;
$X^\mu$ has dimension of length and $\sigma^\alpha$ is dimensionless (it is an angle,
\cref{eq:classical:periodicity}). The integrand of \eqref{eq:classical:ng} then has dimension
length$^2$, so $T$ has dimension length$^{-2}$, which for $c=\hbar=1$ is energy per length, and
$\ap$ is a length squared. For fundamental strings $\ell_s$ is the only dimensionful
parameter of the theory. Its value is not known; the tension is typically taken an order of
magnitude or so below the Planck scale, $T\lesssim(10^{18}\ \text{GeV})^2$
\source{Tong §1.2, ll.~1143--1156}.

\begin{example}[The tension of a QCD flux tube in everyday units]\label{ex:classical:qcd}
The Nambu--Goto action also describes, approximately, string-like objects of finite width
such as the chromo-electric flux tubes of QCD, as long as they are nearly straight on the
scale of their width; for those $T\sim(1\ \text{GeV})^2$ \source{Tong §1.2, ll.~1124--1142}.
Take $T=1\ \text{GeV}^2$ exactly to see the sizes involved. Restoring units with
$\hbar c=0.1973\ \text{GeV\,fm}$ (standard value; not taken from this book's library),
\[
  T=\frac{1\ \text{GeV}^2}{\hbar c}=5.07\ \frac{\text{GeV}}{\text{fm}}
   =5.07\times\frac{1.602\times10^{-10}\ \text{J}}{10^{-15}\ \text{m}}\approx 8\times10^{5}\ \text{N},
\]
the weight of about eighty tonnes, acting between two quarks. The corresponding string length
is $\ell_s=\sqrt{\ap}=(2\pi T)^{-1/2}=0.399\ \text{GeV}^{-1}=0.079$ fm. A fundamental string
with $T$ near $(10^{18}\ \text{GeV})^2$ has a tension larger by a factor $10^{36}$ and a
length $\ell_s$ smaller by a factor $10^{18}$.
\end{example}

\subsection{Symmetries, momenta and equations of motion}\label{sec:classical:ngeom}

The Nambu--Goto action has two symmetries of different nature
\source{Tong §1.2.1, ll.~1159--1170}:
\begin{itemize}[nosep]
\item \emph{Poincar\'e invariance} $X^\mu\to\Lambda^\mu{}_\nu X^\nu+c^\mu$. From the
worldsheet point of view this is a \emph{global} symmetry: $\Lambda$ and $c$ do not depend on
$\sigma^\alpha$. By Noether's theorem it yields conserved momentum and angular momentum.
\item \emph{Reparametrization invariance} $\sigma^\alpha\to\tilde\sigma^\alpha(\sigma)$, a
\emph{gauge} symmetry, shown above.
\end{itemize}
Define the worldsheet momentum densities, with ${\cal L}$ the integrand of
\eqref{eq:classical:ng},
\begin{equation}\label{eq:classical:ngmomenta}
  \Pi^\tau_\mu=\frac{\partial{\cal L}}{\partial\dot X^\mu}
   =-T\,\frac{(\dot X\cdot X')X'_\mu-X'^2\,\dot X_\mu}{\sqrt{(\dot X\cdot X')^2-\dot X^2X'^2}},
  \qquad
  \Pi^\sigma_\mu=\frac{\partial{\cal L}}{\partial X'^\mu}
   =-T\,\frac{(\dot X\cdot X')\dot X_\mu-\dot X^2\,X'_\mu}{\sqrt{(\dot X\cdot X')^2-\dot X^2X'^2}} .
\end{equation}
The Euler--Lagrange equations are
\begin{equation}\label{eq:classical:ngeom}
  \partial_\tau\Pi^\tau_\mu+\partial_\sigma\Pi^\sigma_\mu=0
  \qquad\Longleftrightarrow\qquad
  \partial_\alpha\big(\sqrt{-\det\gamma}\,\gamma^{\alpha\beta}\partial_\beta X^\mu\big)=0
\end{equation}
\source{Tong §1.2.2, ll.~1171--1208}. The second form uses $\delta\sqrt{-\gamma}=\tfrac12
\sqrt{-\gamma}\,\gamma^{\alpha\beta}\delta\gamma_{\alpha\beta}$ and
$\delta\gamma_{\alpha\beta}=2\,\partial_\alpha X\cdot\partial_\beta\delta X$ (by symmetry in
$\alpha\beta$), followed by an integration by parts. Both forms are non-linear in $X$.

As for the particle, the momenta are constrained. Contracting
\eqref{eq:classical:ngmomenta} with $X'^\mu$ and with itself gives, identically,
\begin{equation}\label{eq:classical:ngconstraints}
  \Pi^\tau\cdot X'=0,\qquad (\Pi^\tau)^2+T^2X'^2=0 .
\end{equation}
The first is immediate: the numerator becomes $(\dot X\cdot X')X'^2-X'^2(\dot X\cdot X')$. For
the second, the square of the numerator is $X'^2\big[X'^2\dot X^2-(\dot X\cdot X')^2\big]$,
which is $-X'^2$ times the expression under the root. (Both identities were also confirmed by
a symbolic check for $D=4$.) There are two constraints at each point of the string, as the
counting in the box of \cref{sec:classical:particle} predicted. They are the
string's version of $p^2+m^2=0$, and they will reappear in a simpler gauge as
\eqref{eq:classical:constraints}.

%=====================================================================================
\section{The Polyakov action}\label{sec:classical:polyakov}

\subsection{Trading the square root for a worldsheet metric}

The square root makes the Nambu--Goto action awkward, in particular in a path integral. The
cure is the one that worked for the particle: introduce an independent metric
$g_{\alpha\beta}(\tau,\sigma)$ on the worldsheet, of signature $(-,+)$, and write
\begin{equation}\label{eq:classical:polyakov}
  S_{\rm P}=-\frac{1}{4\pi\ap}\int\dd^2\sigma\,\sqrt{-g}\;g^{\alpha\beta}\,
  \partial_\alpha X^\mu\partial_\beta X^\nu\,\eta_{\mu\nu},\qquad g=\det g_{\alpha\beta},
\end{equation}
the \emph{Polyakov action}\index{Polyakov action} \source{Tong §1.3, ll.~1209--1227}. From
the worldsheet point of view it describes $D$ massless scalar fields $X^\mu$ coupled to
two-dimensional gravity. Note the prefactor: $1/(4\pi\ap)=T/2$.

\begin{historybox}
The action \eqref{eq:classical:polyakov} was not discovered by Polyakov; it carries his name
because he understood how to use it in the path integral
\source{Tong §1.3, ll.~1220--1223}. That use is the subject of \cref{ch:ghosts}.
\end{historybox}

\paragraph{Equation of motion for $X$.} Varying $X^\mu$ with $g_{\alpha\beta}$ fixed gives
\begin{equation}\label{eq:classical:poly-eomX}
  \partial_\alpha\big(\sqrt{-g}\,g^{\alpha\beta}\partial_\beta X^\mu\big)=0,
\end{equation}
the same as \eqref{eq:classical:ngeom} but with $g$ in place of $\gamma$
\source{Tong §1.3, ll.~1228--1235}.

\paragraph{Equation of motion for $g$.} We need $\delta\sqrt{-g}$. From
$\delta\det g=\det g\; g^{\alpha\beta}\delta g_{\alpha\beta}$ and
$g^{\alpha\beta}\delta g_{\alpha\beta}=-g_{\alpha\beta}\delta g^{\alpha\beta}$ (vary
$g^{\alpha\beta}g_{\beta\gamma}=\delta^\alpha_\gamma$),
\begin{equation}\label{eq:classical:varsqrtg}
  \delta\sqrt{-g}=\tfrac12\sqrt{-g}\,g^{\alpha\beta}\delta g_{\alpha\beta}
   =-\tfrac12\sqrt{-g}\,g_{\alpha\beta}\,\delta g^{\alpha\beta}.
\end{equation}
Therefore
\begin{equation}\label{eq:classical:varg}
  \delta S_{\rm P}=-\frac{T}{2}\int\dd^2\sigma\,\sqrt{-g}\;\delta g^{\alpha\beta}
  \Big(\partial_\alpha X\cdot\partial_\beta X-\tfrac12\,g_{\alpha\beta}\,g^{\rho\sigma}\partial_\rho X\cdot\partial_\sigma X\Big)=0
\end{equation}
\source{Tong §1.3, ll.~1236--1256}. The bracket must vanish. Recognizing
$\partial_\alpha X\cdot\partial_\beta X=\gamma_{\alpha\beta}$, this says
\begin{equation}\label{eq:classical:gsol}
  g_{\alpha\beta}=2f(\tau,\sigma)\,\gamma_{\alpha\beta},\qquad
  f^{-1}=g^{\rho\sigma}\gamma_{\rho\sigma}
\end{equation}
\source{Tong §1.3, ll.~1257--1268}: the worldsheet metric equals the induced metric up to a
positive function $f$ that the equation leaves undetermined. (Insert
$g=2f\gamma$ into the definition of $f$: $g^{\rho\sigma}\gamma_{\rho\sigma}=(2f)^{-1}
\gamma^{\rho\sigma}\gamma_{\rho\sigma}=(2f)^{-1}\cdot2=f^{-1}$, an identity; $f>0$ is needed
for $g$ to have the signature of $\gamma$.)

\paragraph{Classical equivalence.} Substitute \eqref{eq:classical:gsol} into
\eqref{eq:classical:polyakov}. In two dimensions $\sqrt{-\det(2f\gamma)}=2f\sqrt{-\det\gamma}$
and $g^{\alpha\beta}=(2f)^{-1}\gamma^{\alpha\beta}$, so
\begin{equation}\label{eq:classical:equiv}
  \sqrt{-g}\,g^{\alpha\beta}\gamma_{\alpha\beta}
  =2f\sqrt{-\gamma}\cdot\frac{1}{2f}\,\gamma^{\alpha\beta}\gamma_{\alpha\beta}=2\sqrt{-\gamma},
  \qquad S_{\rm P}\to-\frac{2}{4\pi\ap}\int\dd^2\sigma\sqrt{-\gamma}=S_{\rm NG}
\end{equation}
\source{Tong §1.3, ll.~1269--1281}. The function $f$ cancels, both here and in
\eqref{eq:classical:poly-eomX}. This cancellation signals a new symmetry.

\subsection{Symmetries: Poincar\'e, diffeomorphisms and Weyl}\label{sec:classical:symm}

The Polyakov action has \source{Tong §1.3.1, ll.~1283--1327}:
\begin{enumerate}[nosep]
\item \emph{Poincar\'e invariance}, global on the worldsheet, as before.
\item \emph{Diffeomorphism invariance}\index{diffeomorphism invariance}
$\sigma^\alpha\to\tilde\sigma^\alpha(\sigma)$, under which $X^\mu$ is a scalar,
$\tilde X^\mu(\tilde\sigma)=X^\mu(\sigma)$, and $g_{\alpha\beta}$ a tensor,
$\tilde g_{\alpha\beta}(\tilde\sigma)=\frac{\partial\sigma^\gamma}{\partial\tilde\sigma^\alpha}
\frac{\partial\sigma^\delta}{\partial\tilde\sigma^\beta}g_{\gamma\delta}(\sigma)$.
Infinitesimally, for $\tilde\sigma^\alpha=\sigma^\alpha-\eta^\alpha(\sigma)$:
$\delta X^\mu=\eta^\alpha\partial_\alpha X^\mu$ and
$\delta g_{\alpha\beta}=\nabla_\alpha\eta_\beta+\nabla_\beta\eta_\alpha$, with $\nabla$ the
Levi-Civita connection of $g$. Compare with the einbein transformations of
\cref{sec:classical:einbein}.
\item \emph{Weyl invariance}\index{Weyl invariance}, which is new:
\begin{equation}\label{eq:classical:weyl}
  X^\mu(\sigma)\to X^\mu(\sigma),\qquad g_{\alpha\beta}(\sigma)\to\Omega^2(\sigma)\,g_{\alpha\beta}(\sigma).
\end{equation}
\end{enumerate}
A Weyl transformation is not a change of coordinates. It is a position-dependent change of
the unit of length on the worldsheet which preserves angles. The Polyakov string regards
two metrics related by \eqref{eq:classical:weyl} as the same physical configuration.

To verify the invariance, and to see why it is special to strings, do the computation on a
worldvolume of dimension $n$. Under \eqref{eq:classical:weyl},
$\det g\to\Omega^{2n}\det g$ and $g^{\alpha\beta}\to\Omega^{-2}g^{\alpha\beta}$, so
\begin{equation}\label{eq:classical:weyl-n}
  \sqrt{-g}\,g^{\alpha\beta}\;\longrightarrow\;\Omega^{\,n-2}\,\sqrt{-g}\,g^{\alpha\beta}.
\end{equation}
The combination that enters \eqref{eq:classical:polyakov} is invariant if and only if $n=2$
\source{Tong §1.3.1, ll.~1328--1338}. A particle ($n=1$) or a membrane ($n=3$) has no such
symmetry. Even for $n=2$ the symmetry is fragile: a potential $\int\dd^2\sigma\sqrt{-g}\,V(X)$
or a worldsheet cosmological constant $\mu\int\dd^2\sigma\sqrt{-g}$ scales with $\Omega^2$ and
breaks it \source{Tong §1.3.1, ll.~1339--1356}. In the quantum theory the requirement of
Weyl invariance becomes much stronger; it fixes the spacetime dimension (\cref{ch:quantum},
\cref{ch:ghosts}) and yields the equations of motion of the background fields
(\cref{ch:backgrounds}).

\begin{pitfallbox}[Common pitfall: three metrics, one symbol]
Three different metrics are in play and sources often write them all as ``$g$'' or
``$\eta$''. (i) The spacetime metric $\eta_{\mu\nu}$, fixed, with indices $\mu,\nu$.
(ii) The induced metric $\gamma_{\alpha\beta}$ of \eqref{eq:classical:induced}, a function of
the $X^\mu$, not an independent field. (iii) The worldsheet metric $g_{\alpha\beta}$ of the
Polyakov action, an independent field. On solutions $g=2f\gamma$, and only the conformal
class of $g$ is determined. When conformal gauge sets $g_{\alpha\beta}=\eta_{\alpha\beta}$
below, that $\eta$ is a two-by-two worldsheet object and has nothing to do with (i); in
particular $\gamma_{\alpha\beta}$ is then \emph{not} equal to $\eta_{\alpha\beta}$ but to
$\eta_{\alpha\beta}/(2f)$.
\end{pitfallbox}

%=====================================================================================
\section{Conformal gauge: wave equation plus constraints}\label{sec:classical:gauge}

\subsection{Counting and fixing}

A symmetric two-by-two metric has three independent components. Diffeomorphisms provide two
arbitrary functions, a Weyl transformation a third. One therefore expects to be able to
remove the metric completely \source{Tong §1.3.2, ll.~1357--1380}. In two steps: first use
the two reparametrizations to bring the metric to the \emph{conformally flat} form
\begin{equation}\label{eq:classical:confgauge}
  g_{\alpha\beta}=\ee^{2\phi(\tau,\sigma)}\,\eta_{\alpha\beta},\qquad \eta_{\alpha\beta}=\diag(-1,+1),
\end{equation}
which is called \emph{conformal gauge}\index{conformal gauge}; then use a Weyl transformation
with $\Omega=\ee^{-\phi}$ to set $\phi=0$, so that $g_{\alpha\beta}=\eta_{\alpha\beta}$. That
the first step is possible, at least locally, is an existence statement for a system of
partial differential equations. A quick argument in Euclidean signature: under
$g'_{\alpha\beta}=\ee^{2\phi}g_{\alpha\beta}$ the Ricci scalar obeys
$\sqrt{g'}R'=\sqrt{g}\,(R-2\nabla^2\phi)$; choose $\phi$ to solve $2\nabla^2\phi=R$, so that
$R'=0$; in two dimensions the Riemann tensor is $R_{\alpha\beta\gamma\delta}=\frac R2(g_{\alpha
\gamma}g_{\beta\delta}-g_{\alpha\delta}g_{\beta\gamma})$, so $R'=0$ implies that $g'$ is flat
\source{Tong §1.3.2, ll.~1381--1411} \TL.

\subsection{Equations of motion and the stress tensor}

With $g_{\alpha\beta}=\eta_{\alpha\beta}$ the Polyakov action becomes the action of $D$ free
massless scalar fields in two dimensions,
\begin{equation}\label{eq:classical:freeaction}
  S=-\frac{1}{4\pi\ap}\int\dd^2\sigma\,\partial_\alpha X\cdot\partial^\alpha X
   =\frac{1}{4\pi\ap}\int\dd\tau\,\dd\sigma\,\big(\dot X^2-X'^2\big),
\end{equation}
and \eqref{eq:classical:poly-eomX} becomes the wave equation\index{wave equation}
\begin{equation}\label{eq:classical:wave}
  \partial_\alpha\partial^\alpha X^\mu=0\qquad\Longleftrightarrow\qquad \ddot X^\mu-X''^\mu=0
\end{equation}
\source{Tong §1.3.2, ll.~1412--1426}. (In fact any conformal gauge gives the same, since
$\sqrt{-g}g^{\alpha\beta}$ does not depend on $\phi$.) The non-linear equations
\eqref{eq:classical:ngeom} have become linear. This cannot be the whole story, and it is not:
fixing $g$ by a gauge choice does not remove its equation of motion. The equation
\eqref{eq:classical:varg} must still be imposed. Define the worldsheet
\emph{stress-energy tensor}\index{stress-energy tensor}
\begin{equation}\label{eq:classical:Tdef}
  T_{\alpha\beta}=-\frac{2}{T}\frac{1}{\sqrt{-g}}\frac{\partial S_{\rm P}}{\partial g^{\alpha\beta}}
  \;\;\xrightarrow{\;g=\eta\;}\;\;
  T_{\alpha\beta}=\partial_\alpha X\cdot\partial_\beta X-\tfrac12\eta_{\alpha\beta}\,\eta^{\rho\sigma}\partial_\rho X\cdot\partial_\sigma X
\end{equation}
\source{Tong §1.3.2, ll.~1427--1447}. (The stress tensor always carries indices; the bare
letter $T$ remains the tension.) The equation of motion of the metric is
$T_{\alpha\beta}=0$. Write out the components with $\eta^{\rho\sigma}\partial_\rho
X\cdot\partial_\sigma X=-\dot X^2+X'^2$:
\begin{align}
  T_{01}&=\dot X\cdot X', \nonumber\\
  T_{00}&=\dot X^2+\tfrac12(-\dot X^2+X'^2)=\tfrac12(\dot X^2+X'^2),\label{eq:classical:Tcomp}\\
  T_{11}&=X'^2-\tfrac12(-\dot X^2+X'^2)=\tfrac12(\dot X^2+X'^2). \nonumber
\end{align}
The trace $\eta^{\alpha\beta}T_{\alpha\beta}=-T_{00}+T_{11}$ vanishes identically. That is a
consequence of Weyl invariance: a Weyl variation is $\delta g^{\alpha\beta}=-2\phi\,
g^{\alpha\beta}$, and invariance of $S_{\rm P}$ under it for arbitrary $\phi(\sigma)$ forces
$g^{\alpha\beta}T_{\alpha\beta}=0$. Only two of the three equations $T_{\alpha\beta}=0$ are
therefore independent:
\begin{equation}\label{eq:classical:constraints}
  \dot X\cdot X'=0,\qquad \dot X^2+X'^2=0
\end{equation}
\source{Tong §1.3.2, ll.~1448--1460}. \emph{The classical string is the free wave equation
\eqref{eq:classical:wave} subject to the two constraints
\eqref{eq:classical:constraints}.}\index{Virasoro constraints} All the non-linearity of the
Nambu--Goto equations has been moved into the constraints.

These are the constraints \eqref{eq:classical:ngconstraints} in a new guise. In conformal
gauge the momentum density is $\Pi^\tau_\mu=\partial{\cal L}/\partial\dot X^\mu=T\dot X_\mu$,
so $\Pi^\tau\cdot X'=0$ is $\dot X\cdot X'=0$ and $(\Pi^\tau)^2+T^2X'^2=0$ is
$\dot X^2+X'^2=0$. Moreover the canonical Hamiltonian of \eqref{eq:classical:freeaction} is
$H=\frac T2\int\dd\sigma(\dot X^2+X'^2)=T\int\dd\sigma\,T_{00}$. It vanishes on the second
constraint, exactly as $H=0$ for the particle in \eqref{eq:classical:Hzero}.

\subsection{What the constraints mean}\label{sec:classical:static}

Conformal gauge leaves some freedom, described in \cref{sec:classical:residual}, which is
enough to reach \emph{static gauge}\index{static gauge} $X^0=t=R\tau$ for a constant $R$. With
$X^\mu=(t,\vec x)$ we have $\dot X^2=-R^2+\dot{\vec x}^{\,2}$ and $X'^2=\vec x^{\,\prime2}$, and the
system becomes
\begin{equation}\label{eq:classical:staticgauge}
  \ddot{\vec x}-\vec x^{\,\prime\prime}=0,\qquad \dot{\vec x}\cdot\vec x^{\,\prime}=0,\qquad
  \dot{\vec x}^{\,2}+\vec x^{\,\prime2}=R^2
\end{equation}
\source{Tong §1.3.2, ll.~1461--1489}. The first constraint says that the velocity is
perpendicular to the string. A motion of the string along itself would only relabel its
points, which is a gauge transformation; the physical oscillations are
\emph{transverse}\index{transverse oscillations}. The second constraint fixes the meaning of
$R$: at an instant when $\dot{\vec x}=0$, $|\vec x^{\,\prime}|=R$, so the length of the string is
$\int_0^{2\pi}\dd\sigma\,|\vec x^{\,\prime}|=2\pi R$ \source{Tong §1.3.2, ll.~1490--1498}. When
the string moves, the same equation trades stretching $|\vec x^{\,\prime}|$ against speed: since
$\dd\vec x/\dd t=\dot{\vec x}/R$, it reads $|\dd\vec x/\dd t|^2+|\vec x^{\,\prime}|^2/R^2=1$. A
point where the string is momentarily unstretched, $\vec x^{\,\prime}=0$, moves at the speed
of light.

\subsection{Residual symmetry}\label{sec:classical:residual}

Introduce worldsheet \emph{lightcone coordinates}\index{lightcone coordinates (worldsheet)}
\begin{equation}\label{eq:classical:lc}
  \sigma^\pm=\tau\pm\sigma,\qquad \partial_\pm=\tfrac12(\partial_\tau\pm\partial_\sigma),\qquad
  \dd s^2=-\dd\tau^2+\dd\sigma^2=-\dd\sigma^+\dd\sigma^- .
\end{equation}
A reparametrization of the special form $\sigma^+\to\tilde\sigma^+(\sigma^+)$,
$\sigma^-\to\tilde\sigma^-(\sigma^-)$ multiplies the flat metric by a function,
$\dd\tilde\sigma^+\dd\tilde\sigma^-=(\partial_+\tilde\sigma^+)(\partial_-\tilde\sigma^-)\,
\dd\sigma^+\dd\sigma^-$, and this factor can be removed by a compensating Weyl transformation.
The combination preserves $g_{\alpha\beta}=\eta_{\alpha\beta}$
\source{Tong §2.2, ll.~1937--1957}. The counting of \cref{sec:classical:gauge} was not wrong:
these transformations depend on functions of one variable only, a negligible subset of the
functions of two variables that we used up. They are the \emph{conformal transformations}%
\index{conformal transformations} of the flat worldsheet and the reason why conformal field
theory (\cref{ch:cft}) is the natural language of the quantum string. For the classical
string they complete the count of degrees of freedom: the general solution of
\eqref{eq:classical:wave} contains $2D$ functions of one variable; the two constraints remove
two; the residual symmetry removes two more; $2(D-2)$ remain, the left- and right-moving
waves in the $D-2$ transverse directions \source{Tong §2.2, ll.~1966--1985}.

%=====================================================================================
\section{Mode expansion: left- and right-movers}\label{sec:classical:modes}

\subsection{General solution}

In lightcone coordinates $\partial_\alpha\partial^\alpha=-\partial_\tau^2+\partial_\sigma^2
=-4\partial_+\partial_-$, so the wave equation is $\partial_+\partial_-X^\mu=0$ and its
general solution is
\begin{equation}\label{eq:classical:LR}
  X^\mu(\tau,\sigma)=X_L^\mu(\sigma^+)+X_R^\mu(\sigma^-)
\end{equation}
with arbitrary functions $X_L$, $X_R$ \source{Tong §1.4, ll.~1500--1512}. A function of
$\sigma^+=\tau+\sigma$ keeps its shape and moves towards decreasing $\sigma$ as $\tau$ grows;
it is called \emph{left-moving}\index{left-movers and right-movers}. A function of $\sigma^-$
is \emph{right-moving}. The two sets of waves pass through each other without interacting;
see \cref{fig:classical:movers}. This decoupling is the single most useful fact about the
closed string, and it survives quantization: the quantum closed string is, to a first
approximation, two copies of one structure, tied together only by the zero modes and by one
constraint (level matching, below).

\begin{figure}[t]
\centering
\begin{tikzpicture}[>=Stealth,scale=0.9]
  \draw[->] (-0.2,0)--(6.8,0) node[below]{$\sigma$};
  \draw[->] (0,-0.2)--(0,3.9) node[left]{$\tau$};
  \draw (6.28,0.08)--(6.28,-0.08) node[below]{\small $2\pi$};
  \draw (0,0.08)--(0,-0.08) node[below]{\small $0$};
  \draw[densely dotted] (6.28,0)--(6.28,3.7);
  % right-mover: sigma^- const, slope +1
  \foreach \s in {0.6,1.0}{\draw[tierL,thick] (\s,0)--(\s+3.2,3.2);}
  \node[tierL,above right] at (3.6,3.2){\small $\sigma^-=\text{const}$};
  % left-mover: sigma^+ const, slope -1
  \foreach \s in {5.2,5.6}{\draw[tierC,thick] (\s,0)--(\s-3.2,3.2);}
  \node[tierC,above left] at (2.6,3.2){\small $\sigma^+=\text{const}$};
  \draw[->,tierL] (1.0,-0.75)--(2.0,-0.75) node[right]{\small right-mover $X_R(\tau-\sigma)$};
  \draw[->,tierC] (5.6,-1.25)--(4.6,-1.25) node[left]{\small left-mover $X_L(\tau+\sigma)$};
\end{tikzpicture}
\caption{The worldsheet cylinder cut open along $\sigma=0\sim2\pi$. A right-moving pulse is
constant along lines $\sigma^-=\tau-\sigma=\text{const}$ and travels towards increasing
$\sigma$ at unit coordinate speed; a left-moving pulse travels the other way. The two pass
through each other unchanged.}
\label{fig:classical:movers}
\end{figure}

Now impose periodicity \eqref{eq:classical:periodicity}. A periodic function is a Fourier
series; but $X_L$ and $X_R$ need not be separately periodic, only their sum. The derivative
$\partial_+X_L$ must be periodic (it equals $\partial_+X$), so $X_L$ is a Fourier series plus
a term linear in $\sigma^+$, and likewise $X_R$. The standard parametrization is
\begin{align}
  X_L^\mu(\sigma^+)&=\tfrac12x^\mu+\tfrac12\ap\,p^\mu\,\sigma^+
     +\ii\sqrt{\frac{\ap}{2}}\sum_{n\neq0}\frac1n\,\tilde\alpha_n^\mu\,\ee^{-\ii n\sigma^+},\nonumber\\
  X_R^\mu(\sigma^-)&=\tfrac12x^\mu+\tfrac12\ap\,p^\mu\,\sigma^-
     +\ii\sqrt{\frac{\ap}{2}}\sum_{n\neq0}\frac1n\,\alpha_n^\mu\,\ee^{-\ii n\sigma^-}
     \label{eq:classical:modeexp}
\end{align}
\source{Tong §1.4, ll.~1513--1538}\index{mode expansion}. The normalizations are chosen for
later convenience. Three remarks make the formula concrete.

\emph{(i) The linear terms.} Their sum is $\tfrac12\ap p^\mu(\sigma^++\sigma^-)=\ap p^\mu\tau$,
which is periodic in $\sigma$ although neither piece is. Had we allowed different
coefficients $p_L^\mu\neq p_R^\mu$, the sum would contain
$\tfrac12\ap(p_L-p_R)^\mu\sigma$, which is not periodic. In flat, non-compact spacetime
periodicity therefore forces equal left and right zero modes. This is exactly the point
that changes in \cref{ch:circle}; see \cref{rem:classical:circle}.

\emph{(ii) Centre of mass.} Averaging over the string,
$\frac{1}{2\pi}\int_0^{2\pi}\dd\sigma\,X^\mu=x^\mu+\ap p^\mu\tau$: the oscillators drop out and
the centre of mass moves on a straight line. The constant $p^\mu$ is the total spacetime
momentum. Indeed, the Noether current of the translation symmetry $X^\mu\to X^\mu+c^\mu$ of
\eqref{eq:classical:freeaction} is $P^\alpha_\mu=-T\,\partial^\alpha X_\mu$, conserved by
\eqref{eq:classical:wave}, and its charge is
\begin{equation}\label{eq:classical:momentum}
  \int_0^{2\pi}\dd\sigma\,P^\tau_\mu=T\int_0^{2\pi}\dd\sigma\,\dot X_\mu
  =\frac{1}{2\pi\ap}\cdot2\pi\cdot\ap\,p_\mu=p_\mu
\end{equation}
\source{Tong §1.4, ll.~1545--1549}. This is why the factor $\ap$ accompanies $p^\mu$ in
\eqref{eq:classical:modeexp}.

\emph{(iii) Reality.} $X^\mu$ is real, hence
\begin{equation}\label{eq:classical:reality}
  \alpha_{-n}^\mu=(\alpha_n^\mu)^\ast,\qquad \tilde\alpha_{-n}^\mu=(\tilde\alpha_n^\mu)^\ast
\end{equation}
\source{Tong §1.4, ll.~1550--1563}. Each pair $(\alpha_n^\mu,\alpha_{-n}^\mu)$, $n>0$, is the
complex amplitude of one right-moving harmonic of frequency $n$. In \cref{ch:quantum} these
amplitudes become annihilation and creation operators.

\subsection{The constraints in modes: \texorpdfstring{$L_n$}{Ln} and \texorpdfstring{$\tilde L_n$}{Ln-tilde}}

From \eqref{eq:classical:lc}, $\partial_\pm X=\tfrac12(\dot X\pm X')$, so
\begin{equation}\label{eq:classical:lcconstraints}
  (\partial_\pm X)^2=\tfrac14\big(\dot X^2+X'^2\pm2\dot X\cdot X'\big)=\tfrac12\big(T_{00}\pm T_{01}\big),
\end{equation}
and the two constraints \eqref{eq:classical:constraints} are equivalent to
$(\partial_+X)^2=(\partial_-X)^2=0$ \source{Tong §1.4.1, ll.~1565--1570}: one constraint on
the left-movers alone, one on the right-movers alone. Differentiate
\eqref{eq:classical:modeexp}; each oscillator term contributes
$\ii\sqrt{\ap/2}\cdot\frac1n\cdot(-\ii n)=\sqrt{\ap/2}$ times $\alpha_n\ee^{-\ii n\sigma^-}$:
\begin{equation}\label{eq:classical:dminusX}
  \partial_-X^\mu=\frac{\ap}{2}p^\mu+\sqrt{\frac{\ap}{2}}\sum_{n\neq0}\alpha_n^\mu\,\ee^{-\ii n\sigma^-}
   =\sqrt{\frac{\ap}{2}}\sum_{n\in\Z}\alpha_n^\mu\,\ee^{-\ii n\sigma^-},
  \qquad \alpha_0^\mu\equiv\sqrt{\frac{\ap}{2}}\,p^\mu .
\end{equation}
Squaring and collecting the coefficient of $\ee^{-\ii n\sigma^-}$,
\begin{equation}\label{eq:classical:Ln}
  (\partial_-X)^2=\frac{\ap}{2}\sum_{m,k}\alpha_m\cdot\alpha_k\,\ee^{-\ii(m+k)\sigma^-}
  =\ap\sum_{n\in\Z}L_n\,\ee^{-\ii n\sigma^-},\qquad
  L_n=\frac12\sum_{m\in\Z}\alpha_{n-m}\cdot\alpha_m ,
\end{equation}
and in the same way $(\partial_+X)^2=\ap\sum_n\tilde L_n\ee^{-\ii n\sigma^+}$ with
$\tilde L_n=\frac12\sum_m\tilde\alpha_{n-m}\cdot\tilde\alpha_m$ and
$\tilde\alpha_0^\mu\equiv\sqrt{\ap/2}\,p^\mu=\alpha_0^\mu$
\source{Tong §1.4.1, ll.~1571--1650}. A Fourier series vanishes only if every coefficient
does. A classical closed string is therefore a set of amplitudes obeying the infinitely many
quadratic conditions
\begin{equation}\label{eq:classical:virasoro}
  L_n=\tilde L_n=0\qquad\text{for all }n\in\Z .
\end{equation}
The $L_n$ are the Fourier modes of the stress tensor. In \cref{ch:cft} they return as
generators of the Virasoro algebra\index{Virasoro generators}; here they are just numbers.

\subsection{The mass formula and level matching}

The conditions with $n=0$ contain $\alpha_0^2=\frac{\ap}{2}p^2$, and $p^2=-M^2$ is minus the
squared rest mass of the string regarded as a particle. Separate the $m=0$ term in $L_0$ and
use $\alpha_{-m}\cdot\alpha_m=\alpha_m\cdot\alpha_{-m}$ (classical amplitudes commute) to
combine $m$ and $-m$:
\begin{equation}\label{eq:classical:L0}
  L_0=\frac12\alpha_0^2+\sum_{n>0}\alpha_{-n}\cdot\alpha_n=-\frac{\ap}{4}M^2+\sum_{n>0}\alpha_{-n}\cdot\alpha_n .
\end{equation}
Setting $L_0=0$ and $\tilde L_0=0$ gives two expressions for the same mass,
\begin{equation}\label{eq:classical:mass}
  M^2=\frac{4}{\ap}\sum_{n>0}\alpha_{-n}\cdot\alpha_n=\frac{4}{\ap}\sum_{n>0}\tilde\alpha_{-n}\cdot\tilde\alpha_n
\end{equation}
\source{Tong §1.4.1, ll.~1651--1679} \TL. The equality of the right-moving and the left-moving
sums is \emph{level matching}\index{level matching}. It holds because
$\alpha_0=\tilde\alpha_0$, which in turn came from periodicity in remark (i). The mass of a
string is the energy stored in its oscillations: by \eqref{eq:classical:reality},
$\alpha_{-n}\cdot\alpha_n=\sum_\mu\eta_{\mu\mu}|\alpha_n^\mu|^2$, a sum of squared
amplitudes (with a minus sign for $\mu=0$, which the constraints $L_{n\neq0}=0$ keep under
control). Classically $M^2$ is any non-negative real number. Quantization
(\cref{ch:quantum}) makes the sums integer-valued, $\sum_{n>0}\alpha_{-n}\cdot\alpha_n\to N$,
and adds a constant from normal ordering, giving the book's convention
$M^2=\frac{2}{\ap}(N+\tilde N-2)$ with $N=\tilde N$ in non-compact flat space.

\begin{remark}[What changes on a circle: a preview]\label{rem:classical:circle}
Suppose one coordinate is periodic, $X\sim X+2\pi R$. Then a closed string may close only up
to this identification, $X(\tau,\sigma+2\pi)=X(\tau,\sigma)+2\pi Rw$ with $w\in\Z$ the
\emph{winding number}\index{winding number}. Remark (i) is then modified: allow
$p_L\neq p_R$ in \eqref{eq:classical:modeexp}. The zero-mode part is
$\tfrac{\ap}{2}(p_L+p_R)\tau+\tfrac{\ap}{2}(p_L-p_R)\sigma$, and matching it to
$\ap p\,\tau+wR\,\sigma$ gives
\begin{equation}\label{eq:classical:pLpR}
  p_L=p+\frac{wR}{\ap},\qquad p_R=p-\frac{wR}{\ap},\qquad
  \alpha_0=\sqrt{\tfrac{\ap}{2}}\,p_R,\quad\tilde\alpha_0=\sqrt{\tfrac{\ap}{2}}\,p_L .
\end{equation}
Now $\alpha_0\neq\tilde\alpha_0$, the two sums in \eqref{eq:classical:mass} differ by
$\frac{\ap}{4}(p_L^2-p_R^2)=p\,wR$, and after quantization, where $p=n/R$, level matching
becomes $N-\tilde N=nw$. This short computation (derived here, not quoted) is the seed of
\cref{ch:circle} and of T-duality in \cref{ch:tduality}: the exchange
$p_R\to-p_R$, $p_L\to p_L$ is the exchange $p\leftrightarrow wR/\ap$.
\end{remark}

%=====================================================================================
\section{Worked example: a pulsating loop}\label{sec:classical:example}

Take $D\geq3$ and look for a circular string in the $(x^1,x^2)$-plane that starts at rest with
radius $R$. In static gauge $X^0=R\tau$ try
\begin{equation}\label{eq:classical:loop}
  \vec x(\tau,\sigma)=R\cos\tau\,(\cos\sigma,\ \sin\sigma).
\end{equation}

\paragraph{Equations of motion and constraints.} $\ddot{\vec x}=-\vec x=\vec x^{\,\prime\prime}$,
so the wave equation holds. Further $\dot{\vec x}=-R\sin\tau(\cos\sigma,\sin\sigma)$ is radial
and $\vec x^{\,\prime}=R\cos\tau(-\sin\sigma,\cos\sigma)$ is tangential, so
$\dot{\vec x}\cdot\vec x^{\,\prime}=0$; and
$\dot{\vec x}^{\,2}+\vec x^{\,\prime2}=R^2\sin^2\tau+R^2\cos^2\tau=R^2$. All of
\eqref{eq:classical:staticgauge} is satisfied. Note that the constant in static gauge had to
be the initial radius: the constraint ties the unit of worldsheet time to the size of the
string.

\paragraph{Motion.} In terms of target-space time $t=R\tau$ the radius is $r(t)=R\,|\cos(t/R)|$
(\cref{fig:classical:loop}). The tension pulls the loop inward; it collapses to a point at
$t=\pi R/2$. Its radial speed is $|\dd r/\dd t|=|\sin(t/R)|$, which reaches the speed of light
at the moment of collapse, in agreement with the remark after
\eqref{eq:classical:staticgauge}, since there $\vec x^{\,\prime}=0$. The loop then re-expands to
radius $R$ at $t=\pi R$. (At the collapse point the Nambu--Goto description is singular, since
$\det\gamma=0$; the conformal-gauge solution passes through it smoothly.)

\begin{figure}[t]
\centering
\begin{tikzpicture}[>=Stealth,scale=0.95]
  \foreach \r/\st in {2.0/thick,1.73/semithick,1.0/thin,0.35/very thin}{\draw[\st] (0,0) circle (\r);}
  \fill (0,0) circle (1.2pt);
  \draw[thin] (35:2.0)--(2.5,1.5) node[right]{\small $\tau=0$};
  \draw[thin] (10:1.73)--(2.5,0.6) node[right]{\small $\tau=\pi/6$};
  \draw[thin] (-20:1.0)--(2.5,-0.4) node[right]{\small $\tau=\pi/3$};
  \foreach \a in {90,180,270}{\draw[->] (\a:1.95)--(\a:1.35);}
  \begin{scope}[xshift=5.6cm,yshift=-1.6cm]
    \draw[->] (0,0)--(5.4,0) node[below]{$t/R$};
    \draw[->] (0,0)--(0,3.4) node[left]{$r/R$};
    \draw[thick,domain=0:1.5708,samples=40] plot ({1.5*\x},{3*cos(\x r)});
    \draw[thick,domain=1.5708:3.1416,samples=40] plot ({1.5*\x},{-3*cos(\x r)});
    \draw (2.356,0.08)--(2.356,-0.08) node[below]{\small $\pi/2$};
    \draw (4.712,0.08)--(4.712,-0.08) node[below]{\small $\pi$};
    \draw (0.08,3)--(-0.08,3) node[left]{\small $1$};
  \end{scope}
\end{tikzpicture}
\caption{The pulsating loop \eqref{eq:classical:loop}. Left: snapshots of the string at
$\tau=0,\pi/6,\pi/3$ and shortly before collapse. Right: radius against time; the slope at
$t=\pi R/2$ is $\mp1$, the speed of light.}
\label{fig:classical:loop}
\end{figure}

\paragraph{Energy and mass.} From \eqref{eq:classical:momentum},
$E=p^0=T\int\dd\sigma\,\dot X^0=2\pi RT=R/\ap$; consistently, comparing $X^0=R\tau$ with the
zero-mode term $\ap p^0\tau$ gives $p^0=R/\ap$. This is the tension times the initial length
$2\pi R$, as it must be for a string initially at rest, and it is conserved although the
length is not: potential energy is converted into kinetic energy. The total spatial momentum
vanishes, so $M=E=R/\ap$.

\paragraph{The same mass from the oscillators.} Use $2\cos\tau\cos\sigma=\cos\sigma^++\cos
\sigma^-$ and $2\cos\tau\sin\sigma=\sin\sigma^+-\sin\sigma^-$ to split
\eqref{eq:classical:loop} into movers:
\[
  X_R^1=\tfrac R2\cos\sigma^-,\quad X_R^2=-\tfrac R2\sin\sigma^-,\qquad
  X_L^1=\tfrac R2\cos\sigma^+,\quad X_L^2=+\tfrac R2\sin\sigma^+ .
\]
Only $n=\pm1$ appear. For $n=\pm1$ the oscillator part of \eqref{eq:classical:modeexp} is
$\ii\sqrt{\ap/2}\,\big(\alpha_1\ee^{-\ii s}-\alpha_1^\ast\ee^{\ii s}\big)
=-2\sqrt{\ap/2}\;\operatorname{Im}\!\big(\alpha_1\ee^{-\ii s}\big)$. With
$c\equiv R/(2\sqrt{2\ap})$ one finds by inspection
\begin{equation}\label{eq:classical:loopmodes}
  \alpha_1^1=-\ii c,\quad \alpha_1^2=-c,\qquad \tilde\alpha_1^1=-\ii c,\quad\tilde\alpha_1^2=+c .
\end{equation}
(Check the first: $\operatorname{Im}(-\ii c\,\ee^{-\ii s})=-c\cos s$, so
$X_R^1=2c\sqrt{\ap/2}\cos s=\frac R2\cos s$.) Then
\[
  \sum_{n>0}\alpha_{-n}\cdot\alpha_n=|\alpha_1^1|^2+|\alpha_1^2|^2=2c^2=\frac{R^2}{4\ap},
  \qquad M^2=\frac4\ap\cdot\frac{R^2}{4\ap}=\frac{R^2}{\ap^2},
\]
in agreement with $M=R/\ap$; the left-movers give the same, so the level-matching condition
holds. The remaining constraints hold too: $L_{\pm1}\propto\alpha_0\cdot\alpha_{\pm1}=0$
because $\alpha_0$ points in the time direction and $\alpha_{\pm1}$ in space, and
$L_2=\tfrac12\alpha_1\cdot\alpha_1=\tfrac12\big((-\ii c)^2+(-c)^2\big)=0$, likewise $\tilde
L_2=\tfrac12\big((-\ii c)^2+c^2\big)=0$. The vanishing of $L_2$ is the statement that the
right-mover is circularly polarized; a linearly polarized mode, say $\alpha_1^2=0$, would
violate it. (All formulas of this example were confirmed by a symbolic check.)

\paragraph{Numbers.} Take $R=10\,\ell_s$. Then $M=10/\ell_s$, the loop collapses after
$t=\pi R/2\approx15.7\,\ell_s$, and the oscillator sum is $R^2/(4\ap)=25$. In the quantum theory
the corresponding level is $N=\tilde N\approx25$ (the normal-ordering constant of
\cref{ch:quantum} shifts the precise value by one). For comparison with
\cref{ch:circle}: a string wound once around a circle of the same radius, $X=R\sigma$, has
$M=2\pi R\,T=R/\ap$ as well, but it cannot collapse, because its winding number is conserved.

%=====================================================================================
\section{What the machine has checked}\label{sec:classical:lean}

\subsection{The honest inventory}

The short answer for this chapter is: nothing about strings. The Lean module assigned to
this chapter, \texttt{StringTheoryFormalization/Foundations/MathlibCore.lean}, contains
exactly one declaration.

\begin{leanbox}[In Lean: the foundation layer]
\begin{leancode}
import Mathlib.Algebra.Category.ModuleCat.Basic
import Mathlib.Analysis.SpecialFunctions.Complex.Circle
import Mathlib.Analysis.InnerProductSpace.Basic
import Mathlib.Topology.Algebra.Module.FiniteDimension
import Mathlib.CategoryTheory.Functor.Basic
import Mathlib.MeasureTheory.Function.L2Space

namespace StringTheory.Foundations

/-- Re-export summary: all Mathlib primitives used by the pipeline.
    This namespace is `open`ed in every downstream block. -/
def mathlib_version : String := "Mathlib4 @ v4.33.1"

end StringTheory.Foundations
\end{leancode}
\textbf{Plain-language reading.} The file imports six Mathlib modules (categories of modules,
the unit circle in $\C$, inner product spaces, finite-dimensional topological vector spaces,
functors, $L^2$ spaces) and defines a constant of type \lean{String}. The definition
\lean{StringTheory.Foundations.mathlib\_version} is a label, not a proposition; it has no
proof because there is nothing to prove. The file's header comment describes it as
``verified'' with no \lean{sorry}; that is true and vacuous. The comment also says that the
file is consumed by all downstream blocks; a text search of the repository finds its name
in eight \texttt{.lean} files including itself. Comments are not checked by the kernel, and
neither is the agreement of the string with the toolchain (the repository's
\texttt{lake-manifest.json} does pin Mathlib at the revision \texttt{v4.33.1}).

\textbf{What it contributes.} Importing a module makes its theorems available. One of them is
relevant to this chapter: Mathlib's \lean{real\_inner\_mul\_inner\_self\_le}, in
\texttt{Mathlib/Analysis/InnerProductSpace/Basic.lean}, is the Cauchy--Schwarz inequality
$\langle u,v\rangle^2\le\langle u,u\rangle\langle v,v\rangle$, the fact that makes the
Euclidean area element \eqref{eq:classical:area} real. It is a theorem of Mathlib, not of
this project.
\end{leanbox}

\subsection{Why an action functional is harder to formalize than its consequences}

This is not an accident of scheduling. Compare what is needed to \emph{state} the two kinds
of result in this chapter.

\emph{The end products are algebra.} The mass formula \eqref{eq:classical:mass}, the zero
modes \eqref{eq:classical:pLpR}, and everything built on them in \cref{ch:circle} and
\cref{ch:tduality} are identities and inequalities among finitely many integers, rationals
or reals. Lean states them in one line, and tactics such as \lean{ring}, \lean{linarith} and
\lean{nlinarith} often prove them in one more. The Tier~A content of this book lives there.

\emph{The starting point is analysis on infinite-dimensional spaces.} To state ``the
Polyakov action is Weyl invariant and its critical points are the solutions of the wave
equation with constraints'' faithfully, one needs at least the following.
\begin{enumerate}[nosep]
\item A type of field configurations: smooth maps $X:\Sigma\to\R^{1,D-1}$ from a
two-manifold, together with Lorentzian metrics on $\Sigma$. These are not Banach spaces, so
the Fr\'echet derivative for normed spaces, which Mathlib has, does not directly apply.
\item The action as an integral over a non-compact $\Sigma$. It diverges for most
configurations; ``stationary'' must be defined through compactly supported variations, and
every integration by parts that we performed silently (from \eqref{eq:classical:varg} or to
\eqref{eq:classical:ngeom}) needs a version of Stokes' theorem with its hypotheses.
\item Lorentzian signature. A text search of the Mathlib revision pinned by this repository
finds Riemannian (positive-definite) metrics on manifolds and lengths of paths, but no file
mentioning pseudo-Riemannian metrics or Euler--Lagrange equations. The expression
$\sqrt{-\det g}$ needs the signature hypothesis merely to be well defined.
\item Gauge redundancy. The physical configuration space is a quotient by diffeomorphisms
and Weyl transformations. ``Every metric can be brought to conformal gauge'' is an existence
theorem for a partial differential equation (isothermal coordinates), true only locally.
\item A dictionary. Even with 1--4 in place, that the resulting mathematical object
\emph{is} a string is a modelling statement, never Tier~A.
\end{enumerate}
None of this is out of reach in principle, but each item is a project of its own, while the
algebra downstream is accessible today. The pointwise algebraic cores of this chapter's
derivations can be isolated and proved: the determinant identity in
\eqref{eq:classical:ng}, the cancellation \eqref{eq:classical:weyl-n}, the einbein
elimination, the tracelessness in \eqref{eq:classical:Tcomp}. Exercises
\ref{exr:classical:leanweyl}--\ref{exr:classical:leangram} ask for them. The authors checked
that the statements given there compile against the repository's Mathlib, but they are not
declarations of the project and carry no tier.

\begin{center}
\small
\begin{tabular}{@{}p{0.50\textwidth}cp{0.36\textwidth}@{}}
\toprule
Claim & Tier & Basis\\
\midrule
Particle: reparametrization invariance, $p^2+m^2=0$, $H=0$ & \TL & Tong §1.1; derived on the page\\
Nambu--Goto action $=$ $-T\times$ area; $T$ is energy per length & \TL & Tong §1.2; derived on the page\\
Constraints \eqref{eq:classical:ngconstraints} & \TL & derived; symbolic check\\
Polyakov $\Leftrightarrow$ Nambu--Goto classically & \TL & Tong §1.3; derived on the page\\
Weyl invariance only for $n=2$ & \TL & Tong §1.3.1; derived on the page\\
Conformal gauge exists locally & \TL & Tong §1.3.2 (sketch only)\\
Mode expansion, $L_n=\tilde L_n=0$, mass formula & \TL & Tong §1.4; derived on the page\\
Pulsating loop, $M=R/\ap$ & \TL & derived; symbolic check\\
$p_{L,R}=p\pm wR/\ap$ (preview) & \TL & derived; see \cref{ch:circle}\\
Cauchy--Schwarz (area is real) & --- & Mathlib theorem, kernel-checked there\\
\lean{mathlib\_version} & \TA & a \lean{String} constant; no mathematical content\\
\bottomrule
\end{tabular}
\end{center}

%=====================================================================================
\begin{summarybox}
\begin{itemize}[nosep,leftmargin=*]
\item The covariant particle action is $-m\times$ proper time. Reparametrization invariance
makes one of the $D$ coordinates unphysical and produces the constraint $p^2+m^2=0$ and
$H=0$. An einbein removes the square root.
\item The Nambu--Goto action is $-T\times$ the area of the worldsheet, computed with the
induced metric $\gamma_{\alpha\beta}=\partial_\alpha X\cdot\partial_\beta X$. $T=1/(2\pi\ap)$
is the energy per unit length, and $\ell_s=\sqrt{\ap}$ is the only scale.
\item The Polyakov action introduces an independent worldsheet metric. Its equation of
motion gives $g=2f\gamma$; substituting back returns Nambu--Goto. It has Poincar\'e,
diffeomorphism and Weyl invariance; the last exists only for two-dimensional worldvolumes.
\item In conformal gauge $g_{\alpha\beta}=\eta_{\alpha\beta}$ the string obeys
$\ddot X=X''$ with constraints $\dot X\cdot X'=0$, $\dot X^2+X'^2=0$, i.e.\ $T_{\alpha\beta}=0$.
Physical oscillations are transverse; $2(D-2)$ functions of one variable remain.
\item $X=X_L(\sigma^+)+X_R(\sigma^-)$. Periodicity forces equal zero modes
$\alpha_0=\tilde\alpha_0=\sqrt{\ap/2}\,p$ in non-compact space. The constraints are
$L_n=\tilde L_n=0$; $n=0$ gives $M^2=\frac4\ap\sum_{n>0}\alpha_{-n}\cdot\alpha_n$ and level
matching.
\item On a circle the zero modes split, $p_{L,R}=p\pm wR/\ap$: the entry point of
\cref{ch:circle}.
\item Tier status: everything here is Tier~L. The assigned Lean file contains a version
string and imports. Formalizing an action principle needs global analysis that is not yet
available in the library; the algebra that follows from it is what later chapters check.
\end{itemize}
\end{summarybox}

%=====================================================================================
\section*{Exercises}
\addcontentsline{toc}{section}{Exercises}

\begin{exercise}[$\star$]
Expand \eqref{eq:classical:particle-frame} for $|\dot{\vec x}|\ll1$ to second order and identify
the two terms. Then show directly from \eqref{eq:classical:particle-p} that $\vec p/E$ is the
velocity.
\end{exercise}

\begin{exercise}[$\star$]
Set $m=0$ in the einbein action \eqref{eq:classical:einbein}. Derive the equations of motion
of $X^\mu$ and of $e$ and show that the particle moves on a null straight line. Verify that
the action is invariant under the transformations $\delta X^\mu=\eta\dot X^\mu$,
$\delta e=\frac{\dd}{\dd\tau}(\eta e)$ up to a total derivative.
\end{exercise}

\begin{exercise}[$\star$]
Verify the two identities \eqref{eq:classical:ngconstraints} from
\eqref{eq:classical:ngmomenta}. Then show that in static gauge with $\dot{\vec x}=0$ the density
$-\Pi^\tau_0$ integrates to $T\times$ length, confirming \eqref{eq:classical:tension} from the
Hamiltonian side.
\end{exercise}

\begin{exercise}[$\star\star$]
Consider a $p$-dimensional extended object with worldvolume dimension $n=p+1$ and the action
$-\frac{T_p}{2}\int\dd^n\sigma\sqrt{-g}\,\big(g^{\alpha\beta}\partial_\alpha X\cdot\partial_\beta
X-(n-2)\big)$. Show that the equation of motion of $g_{\alpha\beta}$ is solved by
$g_{\alpha\beta}=\gamma_{\alpha\beta}$ for $n\neq2$, with no free function, and that
substituting back gives $-T_p\int\dd^n\sigma\sqrt{-\gamma}$. Explain the role of
\eqref{eq:classical:weyl-n} in the difference between $n=2$ and $n\neq2$.
\end{exercise}

\begin{exercise}[$\star\star$]\label{exr:classical:trace}
Show, for an arbitrary worldsheet metric and without using the equations of motion, that
$g^{\alpha\beta}T_{\alpha\beta}=0$ for the tensor in the bracket of \eqref{eq:classical:varg}.
Which property of two dimensions did you use?
\end{exercise}

\begin{exercise}[$\star\star$ A rotating folded string and the Regge slope]
In static gauge $X^0=R\tau$ consider $\vec x=R\cos\sigma\,(\cos\tau,\sin\tau)$. (a) Verify
\eqref{eq:classical:staticgauge}. (b) Describe the motion; which points move at the speed of
light? (c) Compute $E=T\int\dd\sigma\,\dot X^0$ and the angular momentum
$J=T\int\dd\sigma\,(x^1\dot x^2-x^2\dot x^1)$ and show $J=\tfrac12\ap E^2$. This linear
relation between spin and mass squared is the origin of the name ``Regge slope'' for $\ap$.
(d) Find the non-zero $\alpha_n^\mu$, $\tilde\alpha_n^\mu$ and check $L_2=\tilde L_2=0$.
\end{exercise}

\begin{exercise}[$\star\star$]
Repeat the derivation of \cref{rem:classical:circle} in full: starting from
\eqref{eq:classical:modeexp} with independent $p_L$, $p_R$ in one compact direction, compute
$L_0-\tilde L_0$ and show that the classical level-matching condition reads
$\sum_{n>0}\alpha_{-n}\cdot\alpha_n-\sum_{n>0}\tilde\alpha_{-n}\cdot\tilde\alpha_n=p\,wR$. Show
that under $R\to\ap/R$ together with $p\leftrightarrow wR/\ap$ the pair $(p_L,p_R)$ goes to
$(p_L,-p_R)$ and that both $L_0$ and $\tilde L_0$ are unchanged.
\end{exercise}

\begin{exercise}[$\star\star$ Lean]\label{exr:classical:leanweyl}
Prove in Lean~4 with Mathlib the pointwise core of \eqref{eq:classical:weyl-n} for $n=2$: a
Weyl rescaling multiplies the determinant of a symmetric two-by-two matrix by $\Omega^4$.
\begin{leancode}
example (Ω a b c : ℝ) :
    (Ω^2 * a) * (Ω^2 * c) - (Ω^2 * b)^2 = Ω^4 * (a * c - b^2) := by
  sorry
\end{leancode}
Then state and prove the einbein elimination of \cref{sec:classical:einbein}: for $m>0$ and
$s>0$ (standing for $\sqrt{-\dot X^2}$),
$\tfrac12\big((s/m)^{-1}(-s^2)-(s/m)\,m^2\big)=-ms$. Hint: \lean{field\_simp} then \lean{ring}.
\end{exercise}

\begin{exercise}[$\star\star\star$ Lean]\label{exr:classical:leangram}
Let $F$ be a real inner product space. Prove that the Gram determinant of two vectors is
non-negative, so that the Euclidean area element \eqref{eq:classical:area} is real. Use the
Mathlib theorem \lean{real\_inner\_mul\_inner\_self\_le}.
\begin{leancode}
example {F : Type*} [NormedAddCommGroup F] [InnerProductSpace ℝ F]
    (u v : F) :
    0 ≤ inner ℝ u u * inner ℝ v v - inner ℝ u v * inner ℝ u v := by
  sorry
\end{leancode}
Then explain in a paragraph why the analogous statement in Minkowski signature,
$(\dot X\cdot X')^2-\dot X^2X'^2\ge0$, is \emph{not} a theorem about arbitrary pairs of
vectors, and formulate a hypothesis on $\dot X$ and $X'$ under which it holds.
\end{exercise}

\begin{exercise}[$\star\star\star$]
Show that a reparametrization $\sigma^\pm\to\tilde\sigma^\pm(\sigma^\pm)$ maps solutions of
\eqref{eq:classical:wave} and \eqref{eq:classical:constraints} to solutions. Use it to show
that any solution can be brought to static gauge $X^0=R\tau$. (Hint: $X^0$ itself solves the
wave equation, so $X^0=X^0_L(\sigma^+)+X^0_R(\sigma^-)$; choose the new coordinates
proportional to $X^0_L$ and $X^0_R$. What must be assumed about these functions?)
\end{exercise}

\section*{Further reading}
\addcontentsline{toc}{section}{Further reading}
\begin{itemize}[nosep,leftmargin=*]
\item The chapter follows \source{Tong §1, ll.~681--1680} throughout: particle §1.1
(ll.~688--938), Nambu--Goto §1.2 (ll.~945--1208), Polyakov §1.3 (ll.~1209--1498), mode
expansions §1.4 (ll.~1500--1680).
\item Residual conformal symmetry and the counting of $2(D-2)$ transverse functions:
\source{Tong §2.2, ll.~1932--1990}; lightcone gauge, which solves the constraints
explicitly, is taken up in \cref{ch:quantum}.
\item Strings of finite width (flux tubes, cosmic strings) and corrections to the
Nambu--Goto action: \source{Tong §1.2, ll.~1124--1142}.
\item For the state of the formal library: \cref{ch:leanprimer} and \cref{ch:architecture};
for the open problem of formalizing variational principles, \cref{ch:open}.
\end{itemize}
